\section*{अध्याय \ref{ch:g5:biodiversity} --- जैव विविधता}

\begin{solution}{exo:g5:biodiversity:1}
किसी एक ख़ास तरह का \omterm{def:g1:living-or-not:living}{सजीव} --- वे व्यक्ति जो एक ही नमूने पर बने हैं और
आपस में जनन करते हैं। उदाहरण: घरेलू गौरैया, आम गुलबहार, लाल लोमड़ी,
बग़ीचे का घोंघा।
\end{solution}

\begin{solution}{exo:g5:biodiversity:2}
किसी जगह में मौजूद जीवन की विविधता --- सबसे बढ़कर उसकी \omterm{def:g5:biodiversity:species}{प्रजातियों} की
संख्या और उनका हाल। फूलों वाला मैदान एक वर्ग मीटर में \omterm{def:g1:plants-around-us:plant}{पौधों} और \omterm{def:g1:animals-around-us:animal}{जंतुओं}
की दर्जनों \omterm{def:g5:biodiversity:species}{प्रजातियाँ} लिए रहता है: ऊँची \omterm{def:g5:biodiversity:biodiversity}{जैव विविधता}। पक्का अहाता लगभग
कोई नहीं।
\end{solution}

\begin{solution}{exo:g5:biodiversity:3}
लगभग बीस लाख। क्योंकि हर गंभीर अभियान में आज भी बिना नाम वाली
\omterm{def:g5:biodiversity:species}{प्रजातियाँ} मिलती हैं --- ख़ास तौर पर भृंग --- इसलिए सूची से आगे लाखों
और होनी ही चाहिए।
\end{solution}

\begin{solution}{exo:g5:biodiversity:4}
एक ही \omterm{def:g2:where-animals-live:habitat}{आवास} में बहुत-सी \omterm{def:g5:biodiversity:species}{प्रजातियाँ} (मैदान का वर्ग मीटर); अगल-बगल
बहुत-से \omterm{def:g2:where-animals-live:habitat}{आवास} (तालाब, बाड़, मैदान); और एक ही \omterm{def:g5:biodiversity:species}{प्रजाति} के भीतर विविधता
(बग़ीचे के किन्हीं दो घोंघों के \omterm{def:g1:animals-around-us:covering}{खोल} हूबहू एक जैसे नहीं)।
\end{solution}

\begin{solution}{exo:g5:biodiversity:5}
समृद्ध जाल में जिस शिकारी का शिकार ख़त्म हो जाए वह दूसरे शिकार पर जा
सकता है, और हर भूमिका के लिए बदली मौजूद रहती है; झटका सह लिया जाता है।
ग़रीब जाल में हर \omterm{def:g5:biodiversity:species}{प्रजाति} के पास कम विकल्प हैं, इसलिए एक का जाना दूसरों
को भी अपने पीछे खींच लेता है।
\end{solution}

\begin{solution}{exo:g5:biodiversity:6}
मधुमक्खियाँ और मक्खियाँ: फूलों का परागण। गुबरैले: पौध-चीलर खाना।
\omterm{prop:g3:sorting-living-things:groups}{पक्षी}: \omterm{ex:g2:animals-grow-up:butterfly}{इल्लियों} की गश्त। केंचुए: मिट्टी को भुरभुरी रखना।
\end{solution}

\begin{solution}{exo:g5:biodiversity:7}
\omterm{def:g2:where-animals-live:habitat}{आवास} उजाड़ना (सुखाया हुआ तालाब); \omterm{def:g2:caring-for-nature:environment}{पर्यावरण} को ज़हरीला करना (कूड़ा और
रसायन); हद से ज़्यादा उठाना (ज़रूरत से ज़्यादा मछली पकड़ना)। \omterm{prop:g5:biodiversity:loss}{विलुप्त} का
अर्थ है पूरी तरह चला जाना --- हर जगह से, हमेशा के लिए।
\end{solution}

\begin{solution}{exo:g5:biodiversity:8}
उत्तर खुला है। घास वाली या लगाई हुई वर्गाकार जगह में आम तौर पर दस या
उससे ज़्यादा \omterm{def:g5:biodiversity:species}{प्रजातियाँ} निकलती हैं; पक्की जगह में लगभग शून्य --- यही
फ़र्क़ वह माप है।
\end{solution}

\begin{solution}{exo:g5:biodiversity:9}
हर \omterm{def:g5:biodiversity:species}{प्रजाति} एक धागा है: उसका जाना हर पड़ोसी के पास --- जो भी उसे खाता
था, और जिसका भी वह परागण या सफ़ाई करती थी --- कम विकल्प छोड़ जाता है,
और जाल की अगला झटका सहने की ताक़त घटा देता है। और चूँकि कोई नहीं जानता
कि एक \omterm{def:g5:biodiversity:species}{प्रजाति} कितने धागे थामे है, इसलिए कोई पहले से यह प्रमाणित नहीं
कर सकता कि उसका जाना हानिरहित है; और विलुप्ति हमेशा के लिए होती है,
इसलिए वह ग़लती कभी सुधारी नहीं जा सकती।
\end{solution}

\begin{solution}{exo:g5:biodiversity:10}
सारसों की एक-एक करके रक्षा का मतलब होता \omterm{prop:g3:sorting-living-things:groups}{पक्षियों} पर पहरा देना; गीले
मैदानों को बहाल करने ने वह फिर से खड़ा किया जिस पर सारस
\emph{निर्भर} हैं --- और \omterm{def:g2:where-animals-live:habitat}{आवास} अपने सारे बाशिंदों का एक साथ पोषण करता
है, इसीलिए सारसों के साथ मेंढक, व्याध-पतंगे और मैदानी \omterm{def:g3:flowers-fruits-seeds:parts}{फूल} भी लौटे।
\omterm{def:g2:where-animals-live:habitat}{आवास} नहीं, तो बाशिंदे नहीं; \omterm{def:g2:where-animals-live:habitat}{आवास} बहाल, तो बाशिंदे वापस।
\end{solution}

\begin{solution}{exo:g5:biodiversity:11}
जो हटाया गया वह सजावट नहीं, एक कार्यबल था: बाड़ और फूलों की \omterm{def:g5:biodiversity:species}{प्रजातियाँ}
परागण करने वालों और \omterm{prop:g3:sorting-living-things:groups}{कीट} खाने वालों को बसाए रखती थीं, और मुफ़्त होने तक
उनकी सेवाएँ किसी के हिसाब में नहीं आती थीं। प्रकृति का काम अदृश्य
बहीखाता है --- लगातार, बिना मज़दूरी, बिना ध्यान --- इसलिए उसका दाम
ठीक तभी पता चलता है जब मज़दूर बेदख़ल कर दिए जाते हैं और वही काम बिल
बनकर लौटते हैं।
\end{solution}
